\chapter{Transformer架构}

\begin{intro}
\begin{itemize}
    \item Transformer架构的工作流程
    \item Transformer架构为什么有用的一些众说纷纭的解释
    \item Transformer架构在NLP和其他领域的应用
\end{itemize}
本章难度较大，部分内容极为艰深。仅阅读如此科普材料是不能深入理解Transformer的。如有兴趣，建议阅读本章引用的许多原始论文。
\end{intro}

十年前，RNN的梯度消失是不可避免的深渊。这时候，Transformer架构横空出世，其论文标题 \textsc{Attention is All You Need} 也更是“不像个论文标题”，但却一语道破了Transformer的核心思想：注意力机制。Transformer架构的出现，彻底颠覆了自然语言处理领域的研究方向，并立刻被应用到其他领域，如计算机视觉、语音识别、图像生成等等，成果斐然，几乎成为了现代人工智能技术的标配。

经过十年的研究，我们已经能够准确地描述Transformer架构的工作流程，但迄今为止“Transformer架构为什么这么有用”“为什么比RNN好那么多”这些问题，还\textbf{没有任何广为接受的的解释}。

换言之，时至今日，Transformer仍然是个可解释性低的黑箱。

本章将带着大家一起走一遍Transformer架构的工作流程，在不涉及“为什么”的前提下，尽可能清晰地描述Transformer 架构是怎样工作的。等你能够清晰地描述了它是怎样工作的之后，可能未来某一天你就是那个能够解释它为什么有效的人了。

\section{Transformer架构的工作流程}

Transformer架构的核心是注意力机制。围绕这个核心，Transformer架构的工作流程可以分为以下几个步骤：
\begin{enumerate}
    \item 输入数据的预处理：将输入数据转换为适合Transformer架构处理的格式，如文本序列、图像块序列等。
    \item 编码器的前向传播：将输入数据通过编码器进行处理，得到编码器的输出表示。编码器一般是多头自注意力层和前馈神经网络的堆叠。
    \item 解码器的前向传播：将编码器的输出表示和解码器的输入数据通过解码器进行处理，得到解码器的输出表示。解码器一般是多头自注意力层、编码器-解码器注意力层和前馈神经网络的堆叠。
    \item 输出数据的后处理：将解码器的输出表示转换为适合任务需求的格式，如文本序列、图像块序列等。
\end{enumerate}

\subsection{单头自注意力机制}

注意力分三部分：查询（Query）、键（Key）、值（Value）。在单头自注意力机制中，输入数据首先被映射为查询、键、值三个向量，然后通过计算查询和键的相似度来得到注意力权重，最后将注意力权重应用于值向量，得到输出表示。

Attention干的事情就是：
\begin{itemize}
    \item 每一个Q去问所有的K：“你们谁和我最相似？”，然后得到一个注意力分布。
    \item 然后用这个注意力分布去加权所有的V，得到一个新的表示。
\end{itemize}
这三个矩阵的建构方法是在模型训练的过程中学出来的。模型本身包含三个权重矩阵$W_Q, W_K, W_V$，分别用于将输入数据映射为查询、键、值三个向量。

在模型的评估（使用）过程中，我们先拿到输入序列$X$，然后算出三个矩阵：
\begin{equation}
    Q = XW_Q, \quad K = XW_K, \quad V = XW_V
\end{equation}
然后计算注意力权重：
\begin{equation}
    \text{Attention}(Q, K, V) = \text{softmax}\left(\frac{QK^\top}{\sqrt{d_k}}\right)V
\end{equation}
其中，$d_k$是键向量的维度，用于对注意力权重进行缩放，将方差控制到1，从而避免出现梯度消失的问题，具体分析见\nameref{sec:scaling_factor}。

为了更好的理解Attention的变化，我们可以假设输入的表示矩阵$X$的形状是$(\text{len}, d_{model})$，其中$\text{len}$是输入序列的长度，$d_{model}$是输入表示的维度。则查询、键、值矩阵的形状分别为：
\begin{itemize}
    \item $Q = XW_Q$，形状为$(\text{len}, d_k)$
    \item $K = XW_K$，形状为$(\text{len}, d_k)$
    \item $V = XW_V$，形状为$(\text{len}, d_v)$
    \item $QK^\top$，形状为$(\text{len}, \text{len})$，这是一个方阵。
    \item 最后注意力矩阵的形状是$(\text{len}, d_v)$，与输入的表示矩阵$X$的形状相同。
\end{itemize}

上述过程中，$d_v$是自己的，而不是Q、K一样的$d_k$。这是因为，在Transformer架构中，Q、K需要相乘，所以不得不一样；但V不需要和Q、K相乘，所以可以不一样。更重要的是，V承载了最重要的信息，$d_v$的大小直接影响了模型的表达能力，所以往往需要单独设置；但$d_v$调大之后会导致计算量的增加，所以需要在表达能力和计算量之间进行权衡。

上述过程中，Q的来源和K、V的来源不一样，且该差异在编码器和解码器中呈现出不同的模式。对于编码器而言，Q、K、V都来自同一个输入序列，即源语言句子（例如“Je suis étudiant”），因此编码器的自注意力是双向的：每个位置都能看到整个序列的所有位置，没有任何屏蔽。但对于解码器，情况更为复杂——原始Transformer解码器的每一层实际包含两个注意力子层\cite{vaswani2023attentionneed}。

第一个子层是掩码自注意力（masked self-attention），也称因果自注意力。其Q、K、V均来自解码器自身的输入序列（即已经生成的目标语言token，如“I am a”），因此本质上也是自注意力。但为了维持自回归生成的性质，必须防止当前位置“偷看”未来的目标词。具体做法是：在计算softmax之前，将$QK^\top$矩阵中所有“查询位置$i$到键位置$j$且$j>i$”的项加上$-\infty$（或一个极大的负数），使得softmax后的权重在这些未来位置上为0。这一操作称为因果掩码（causal mask）\cite{vaswani2023attentionneed}。加掩码后的注意力分布仅覆盖已生成的历史token，保证了预测当前token时模型只能依赖于左侧的上下文。

第二个子层是交叉注意力（cross-attention），在原始论文中称为“编码器-解码器注意力”。其Q来自解码器前一层（掩码自注意力后）的输出，而K和V则来自编码器顶层的输出序列\cite{vaswani2023attentionneed}。这意味着在交叉注意力中，查询向量携带了“当前解码位置需要什么信息”的信号，而键矩阵和值矩阵则编码了整个源语言句子的上下文表示。通过计算Q与K的相似度，解码器能够动态地从源句中提取最相关的信息，再通过加权聚合V获得最终的上下文向量。交叉注意力不再使用因果掩码，因为解码器生成时可以完整地“阅读”整个源语言序列。

总结起来，可以简单记忆为：“\textbf{编码器看自己，解码器自注意力看自己的过去，交叉注意力看编码器的全部}。”为了更精确地反映这种结构，表\cref{tab:QKV}给出了一个具体翻译例子中各注意力子层的Q、K、V来源。
\begin{table}
    \centering\small
    \begin{tabular}{c|ccc|c}
        \toprule
        组件 & Q & K \& V & 序列长度 & 掩码 \\
        \midrule
        编码器自注意力 & Je suis étudiant & Je suis étudiant & 3 & 无（双向）\\
        解码器掩码自注意力 & I am a & I am a & 3 & 因果掩码（只看过去）\\
        解码器交叉注意力 & 解码器输出 & 编码器输出 & Q长3，K/V长3 & 无\\
        \bottomrule
    \end{tabular}
    \caption{原始Transformer中编码器和解码器各注意力子层的Q、K、V来源与掩码情况}
    \label{tab:QKV}
\end{table}

\subsection{多头注意力}

多头注意力的定义是：将查询、键、值分别线性变换为多个子空间中的表示，然后在每个子空间中计算注意力，最后将这些注意力的输出拼接起来，再进行线性变换得到最终的输出。形式化地说，设共有 $h$ 个头，每个头的维度为 $d_k = d_{\text{model}} / h$，则：
\begin{align}
    \text{head}_i &= \text{Attention}(X W_Q^{(i)}, X W_K^{(i)}, X W_V^{(i)}), \quad i = 1, \dots, h, \\
    \text{MultiHead}(X) &= \text{Concat}(\text{head}_1, \dots, \text{head}_h) \, W_O,
\end{align}
其中 $W_Q^{(i)}, W_K^{(i)}, W_V^{(i)} \in \mathbb{R}^{d_{\text{model}} \times d_k}$，$W_O \in \mathbb{R}^{h d_k \times d_{\text{model}}}$。这一设计最早由 Vaswani 等人在原始 Transformer 论文中提出\cite{vaswani2023attentionneed}，并迅速成为现代深度学习中几乎所有序列建模任务的标准组件。

在实际的试验中，单头自注意力只能捕捉到一种关系，而多头中，不同的头有显著的不同偏好：有的更喜欢捕捉局部关系，有的更喜欢捕捉全局关系，有的更喜欢捕捉语义关系，有的更喜欢捕捉语法关系等等。而且目前发现最好的头的数目是 8 或 16，过多或过少都会导致性能下降。

为什么最优的头数通常落在 8 或 16？这个问题的完整答案至今未知，但已有若干线索。一个操作性解释是：在原始 Transformer 的 $d_{\text{model}} = 512$ 配置下，$h = 8$ 意味着每个头的维度是 $d_k = 64$。如果头数太多（如 $h = 32$，则 $d_k = 16$），每个头的子空间维度太低，可能不足以形成有意义的注意力模式；如果头数太少（如 $h = 2$，则 $d_k = 256$），则回到了单头表达能力受限的问题。Michel 等人\cite{michel2019sixteenheadsreallybetter} 的实验表明，仅保留少数几个最重要的头（甚至可以只剩 1--4 个头）就足以在测试时维持大部分性能，但在训练时，多头提供了更丰富的梯度信号和更稳定的优化路径。因此，8 或 16 更可能是训练效率和推理效率之间的一个经验上的优解，而非表达能力的硬性上限。

进一步的探针试验表明，在模型训练到一半的时候把一些头初始化，继续训练模型依然能够收敛，这些被初始化的头依然能够学到东西（但可能和先前学的东西不一样），这说明某个头喜欢什么东西并不是设计使然，而是在训练过程中自然形成的\cite{voita2019analyzingmultiheadselfattentionspecialized,michel2019sixteenheadsreallybetter}。这也与“彩票假说”\cite{frankle2019lotterytickethypothesisfinding} 在注意力头上的延伸研究一致：模型有足够的冗余来容忍头级别的参数扰动。

更大胆的消融试验表明，即使在模型训练到一半的时候一口气删了一半的头，模型的 loss 也仅仅是缓慢上升\cite{michel2019sixteenheadsreallybetter}。Michel等人按重要性（基于掩码训练后对损失的影响）对头进行排序，然后从低重要性的头开始逐个移除，发现在移除大量头后模型性能才出现显著退化。进一步发现，移除头后的 loss 上升是一个缓慢而平滑的过程，说明不同头之间存在显著的\textbf{功能冗余}：多个头可能学到了相似但略有差异的特征，删除其中一些可以由剩余的头部分补偿。

Kobayashi 等人\cite{kobayashi2020attentionweightanalyzingtransformers} 对 BERT 的注意力权重进行了深入分析，发现仅靠注意力权重（$QK^\top$ 的 softmax）并不足以解释模型的行为，Value 向量所携带的\textbf{词汇语义信息}才是最终输出的决定性因素。冻结 $V$ 等价于切断了每个 token 将其自身信息（通过 $W_V$ 投影后的内容表示）传递给其他 token 的路径，使得注意力机制退化为仅对一组固定的、不可学习的“语义基向量”进行加权组合，这自然导致模型表达能力的崩溃。这与前文单头注意力一节中“V 承载了最重要的信息，$d_v$ 的大小直接影响了模型的表达能力”的论断完全一致，说明了多头注意力机制能起作用是很多因素共同作用的结果，而不是单纯的多头设计使然。

综合上述实验，我们可以形成这样一幅图景：多头注意力的鲁棒性源于不同头之间的功能冗余和互补。当部分头被移除或重置时，剩余的头可以通过调整自身的注意力模式进行部分补偿，因此 loss 仅缓慢上升。但 V（Value 投影矩阵及其输出）承载了每个 token 的\textbf{内容信息本身}，它决定了“当某个 token 被认为重要时，什么样的信息会被传递出去”。如果 V 被冻结，信息流动的核心管道被卡死，无论注意力权重如何调整都无法补救，因此 loss 会剧烈爆炸。这再次印证了一个根本的设计智慧：在 Q/K/V 三元组中，Q 和 K 负责\textbf{路由}（决定信息流向），而 V 负责\textbf{内容}（决定传递什么信息）。路由可以冗余、可以出错后修正，但内容一旦断裂，整个系统即告崩溃。

\subsection{自注意力与FNN、RNN的本质区别}

在本讲义的前几章中，我们分别讨论了FNN、RNN。现在我们可以将它们并列对比，以更清晰地理解自注意力的独特位置。

\begin{description}
    \item[全连接层（线性层）] 每个输出神经元与\textbf{所有输入}神经元相连，权重矩阵 $W$ 的每个元素都是独立参数。它的聚合方式是\textbf{全局但固定}的：一旦训练完毕，无论输入什么，$x_i$ 对 $y_j$ 的贡献都由 $W_{ji}$ 固定。参数量为 $O(n_{\text{in}} \cdot n_{\text{out}})$，且不具备平移等变性。
    \item[卷积层] 利用\textbf{局部感受野}和\textbf{权重共享}实现平移等变性。每个输出像素仅依赖于输入的一个小邻域，同一卷积核在整个输入上滑动，参数被所有位置共享。参数量极小（如 $3\times3\times C_{\text{in}}\times C_{\text{out}}$），但感受野受限，需要堆叠多层才能捕捉全局依赖。
    \item[循环层（RNN/LSTM）] 通过隐藏状态将历史信息\textbf{按时间顺序逐步传递}。理论上可以建模任意长的依赖，但实际中受梯度消失/爆炸的限制，有效记忆长度有限（通常不超过几十个时间步）。此外，RNN 的计算是\textbf{串行}的，无法充分利用并行硬件。
    \item[自注意力层] 在单个操作内，\textbf{每个位置都直接关注所有位置}。聚合权重是\textbf{数据依赖}的，不同的输入会产生不同的注意力分布。路径长度为 $O(1)$，不存在时间上的串行瓶颈。代价是 $O(n^2)$ 的计算和存储复杂度。
\end{description}

事实上，现代最优秀的架构（如 Vision Transformer \cite{dosovitskiy2021imageworth16x16words}、Conformer \cite{gulati2020conformerconvolutionaugmentedtransformerspeech}）正是将自注意力与卷积结合起来，取长补短。卷积擅长捕捉局部纹理和边缘，自注意力擅长建模全局语义关联。理解它们各自的设计哲学和适用场景，是设计高效深度学习系统的关键。

\section{为什么自注意力机制如此有效？}

我们已经从直觉上理解了注意力机制“做什么”：让每个查询去问所有的键，然后按相似度聚合值。

但许多重要的问题值得追问：\textbf{为什么这种设计是合理的，又为什么它能够在如此多样的任务上都表现出色？}这些问题困扰了研究者们十年。一篇又一篇关于 Transformer 的论文都在回答问题、发现问题。

\subsection{自注意力机制的合理性}

\paragraph{从固定权重聚合到数据依赖的聚合}

在自注意力机制出现之前，处理序列数据的主流方式是递归神经网络（RNN）和卷积神经网络（CNN）。RNN 通过隐藏状态 $h_t$ 将过去的信息“裹挟”到当前时刻，其聚合方式由参数 $W$ 和 $U$ 固定。无论当前输入是什么，$h_{t-1}$ 对 $h_t$ 的贡献比例都是由相同的权重矩阵决定的。类似地，CNN 的卷积核一旦训练好，对所有输入位置都使用完全相同的权重进行聚合。这在某些场景下是高效且合理的，但它也意味着：\textbf{模型无法根据输入内容动态调整“谁更应该被关注”}。

自注意力机制的核心创新恰恰在于\textbf{数据依赖的聚合}：对于同一个序列，不同的输入（不同的 $Q$）会产生不同的注意力分布。模型的学习目标不再是“某个位置的权重应该是多少”，而是“如何从数据中\textbf{学会判断谁和谁相关}”。正如 Vaswani 等人\cite{vaswani2023attentionneed} 指出的，这种机制使得每一层都可以直接访问序列中所有位置的信息，路径长度为 $O(1)$，而不像 RNN 那样需要 $O(T)$ 的时间步才能建立远距离依赖。

\paragraph{软注意力的归纳偏置：对齐与检索}

从更抽象的角度看，自注意力实现了一种\textbf{软对齐}（soft alignment）操作。Bahdanau 等人\cite{bahdanau2016neuralmachinetranslationjointly} 最早将注意力引入序列到序列模型时，正是为了解决机器翻译中“源语言和目标语言之间的词对齐”问题。注意力机制通过计算 Query 和 Key 之间的相似度，隐含地回答了“当前需要关注输入中的哪些部分”这一问题。与传统的硬对齐（如隐马尔可夫模型中的离散隐状态）不同，软注意力允许模型同时关注多个位置，且关注程度由可微的 softmax 函数平滑控制，因此可以直接通过梯度下降进行端到端优化。

这种“软对齐”视角还可以被进一步推广：如果将 Key 理解为某种“索引”或“地址”，将 Value 理解为存储在该地址的“内容”，那么自注意力本质上就是一次字典查找。模型通过 $QK^\top/\sqrt{d_k}$ 来度量“查询与各条目之间的匹配程度”，再通过 softmax 归一化为概率分布，最终用该分布对 Value 进行加权平均。正因为这一过程完全可微，$W_Q, W_K, W_V$ 都可以通过反向传播学习，使模型能够自动发现序列内部最有用的关联模式。

\paragraph{为什么需要 $Q$、$K$、$V$ 三个独立的投影？}

初学者常问的一个问题是：既然在自注意力中 $Q$、$K$、$V$ 都来自同一个 $X$，为什么需要三个不同的权重矩阵，而不是直接使用 $X$ 自身来做内积和聚合？答案可以分两层来理解。

第一层是\textbf{表达能力}。如果直接用 $XX^\top$ 计算相似度，那么这个相似度矩阵是完全由 $X$ 的各行之间的内积决定的，它只能捕捉一种固定的相似性定义。而引入 $W_Q$ 和 $W_K$ 后，$QK^\top = XW_Q W_K^\top X^\top$ 可以学习到输入在不同子空间中的相似性：例如，一个子空间可能关注“语法角色”的相似性，另一个子空间可能关注“语义内容”的相似性。这是多头注意力的基础，但即使在单头中，$W_Q$ 和 $W_K$ 的分离已经提供了显著更丰富的相似性度量。

第二层是\textbf{非对称性}。相似度矩阵 $QK^\top$ 在 softmax 之后并非对称的：即使 $Q_i$ 和 $K_j$ 的内积等于 $Q_j$ 和 $K_i$ 的内积（当 $W_Q = W_K$ 时确实如此），经过按行 softmax 归一化后，位置 $i$ 对 $j$ 的关注程度与 $j$ 对 $i$ 的关注程度一般是不等的。这种\textbf{方向性的注意}在自然语言中非常重要——主语对谓语的关注程度与谓语对主语的关注程度不同，而这正是语言理解所需要的。$W_Q$ 和 $W_K$ 的分离（以及 $W_V$ 的进一步分离）使模型有能力学习这种非对称的依赖关系\cite{shaw2018selfattentionrelativepositionrepresentations}。

\subsection{关于缩放因子 $\sqrt{d_k}$}\label{sec:scaling_factor}

\paragraph{从方差角度理解缩放的必要性}

假设 Query 向量 $q$ 和 Key 向量 $k$ 的各分量是独立的随机变量，均值为 $0$，方差为 $1$。那么内积 $q^\top k = \sum_{i=1}^{d_k} q_i k_i$ 的均值为 $0$，方差为 $d_k$（因为 $\text{Var}(q_i k_i) = \text{Var}(q_i)\text{Var}(k_i) = 1$，且各乘积项独立）。当 $d_k$ 较大（例如原始 Transformer 中 $d_k = 64$）时，未经缩放的 $q^\top k$ 会具有方差 $64$，对应的标准差约为 $8$。这意味着某些内积值可能达到 $\pm 15$ 甚至 $\pm 20$ 的量级。

当这样量级的数值被送入 softmax 函数时会发生什么？softmax 是 $\exp$ 的归一化：
\[
\text{softmax}(z)_i = \frac{\exp(z_i)}{\sum_j \exp(z_j)}.
\]
当 $z_i$ 的绝对值很大时，$\exp(z_i)$ 会以指数方式放大差异，导致 softmax 的输出极度集中于最大值对应的位置，其余位置的梯度趋近于零。Vaswani 等人\cite{vaswani2023attentionneed} 在论文中指出：“我们期望对于大的$d_k$，内积的值会变得很大，从而将 softmax 函数推入梯度极小的区域。” 这一现象与 sigmoid 的饱和区完全类似：当输入绝对值太大，激活函数的导数趋于零，反向传播中的梯度信号就被“挤死”了。

通过除以 $\sqrt{d_k}$，内积的方差被归一化为 $1$，从而将 softmax 的输入控制在一个合理的数值范围内，避免了梯度消失。可以证明，在温和的假设下（$q$ 和 $k$ 的分量独立且均值为零、方差为 $1$），$\frac{q^\top k}{\sqrt{d_k}}$ 渐进服从标准正态分布 $\mathcal{N}(0, 1)$，其 softmax 输出的熵保持在中等水平，既不过于“尖锐”（导致梯度消失），也不过于“平坦”（导致信息混合）\cite{ba2016layernormalization}。

\begin{example}{缩放对 softmax 分布的定量影响}{}
考虑一个长度为 $4$ 的序列，某次前向传播中，某个 Query 与四个 Key 的内积（未经缩放）为 $[20, 5, -3, 8]$。分别计算未缩放和缩放后的 softmax：

未缩放：$\exp(20) \approx 4.85\times 10^8$，其余 $\exp$ 值约 $1.48\times 10^2$、$4.98\times 10^{-2}$、$2.98\times 10^3$。归一化后，第一个位置的注意力权重几乎为 $1$（约 $0.999$），其余位置接近 $0$。此时的梯度几乎为零，模型丧失了从其他位置学习的能力。

若使用 $d_k=64$，缩放后方差归一化为 $1$，内积变为 $[2.5, 0.625, -0.375, 1.0]$。此时 softmax 输出约为 $[0.72, 0.11, 0.04, 0.13]$，分布较为平滑，梯度可以有效地流向所有位置。
\end{example}

\paragraph{缩放因子与训练稳定性的理论联系}

缩放因子对训练稳定性的贡献不仅体现在前向传播中保持 softmax 的“柔和”，也体现在反向传播中保持梯度的合理量级。自注意力机制中，若内积的方差过大，softmax 的输入会非常尖锐，输出趋近于 one-hot 分布。此时 softmax 的 Jacobian 几乎为零，反向传播中流经该注意力头的梯度将消失殆尽。换句话说，随着方差的增大，我们是从梯度健康走向梯度消失。

因此，缩放因子对训练稳定性的贡献，本质在于保护梯度。若内积方差太大，softmax 进入饱和区，相应的注意力头的梯度几乎为零，该头将无法有效学习。同时，softmax 输出过于集中也会使模型只注意到极少的位置，丧失从其他位置获取信息的能力。因此，$1/\sqrt{d_k}$ 将方差标准化，使得 softmax 的输入大致保持在 $[-3, 3]$ 的范围，从而兼顾了梯度有效性和信息混合的充分性。这也是原始论文中将其引入的根本动机。

值得指出的是，Xiong 等人\cite{xiong2020layernormalizationtransformerarchitecture} 在分析 Transformer 训练的稳定性时发现，即使在有层归一化的条件下，注意力层的梯度仍然可能在深层传播中出现异常。正确设置缩放因子并配合适当的初始化策略（如 Xavier 初始化或 He 初始化），是保证深层 Transformer 可训练的关键因素。

\subsection{注意力矩阵与 Scaling Law}

自注意力机制的计算复杂度和存储开销与序列长度 $n$ 的关系，是评估其可扩展性的核心指标。设序列长度为 $n$，表示维度为 $d$，则：
\begin{itemize}
    \item $QK^\top$ 的计算复杂度为 $O(n^2 d)$，产生的注意力矩阵形状为 $n\times n$。
    \item softmax 逐行归一化的复杂度为 $O(n^2)$。
    \item 注意力矩阵与 $V$ 相乘的复杂度为 $O(n^2 d)$。
\end{itemize}
因此，\textbf{单头自注意力的总时间和空间复杂度均为 $O(n^2 d)$}。对于 $n$ 达到数千甚至上万的“长上下文”场景（例如处理整篇论文、长视频序列或基因组数据），$O(n^2)$ 的存储和计算开销成为严重瓶颈。

这一问题催生了大量的高效注意力变体。例如，Beltagy 等人\cite{beltagy2020longformerlongdocumenttransformer} 提出的 Longformer 将全局注意力与局部滑窗注意力相结合，将复杂度降至 $O(n)$；Kitaev 等人\cite{kitaev2020reformerefficienttransformer} 的 Reformer 使用局部敏感哈希（LSH）将注意力限制在近似近邻内，复杂度降为 $O(n\log n)$；Wang 等人\cite{wang2020linformerselfattentionlinearcomplexity} 的 Linformer 通过理论分析证明自注意力矩阵是低秩的，因而可以将其投影到低维空间进行计算，将复杂度降为 $O(n)$。这些工作从不同角度印证了一个共同的洞察：对于大多数实际序列，\textbf{注意力矩阵的有效秩远小于 $n$}，长距离的精确注意力往往是不必要的，可以通过近似手段来大幅减少计算量。

有趣的是，近年来的 Scaling Law 研究发现\cite{kaplan2020scalinglawsneurallanguage}，增加序列长度 $n$ 对降低模型损失（即提升性能）的贡献大致服从幂律关系 $\text{Loss} \propto n^{-\alpha}$，其中 $\alpha$ 与模型大小和数据量相关。这意味着在某些场景下，\textbf{增加注意力窗口确实比增加模型参数更“划算”}。这也是为什么现代大语言模型竞相扩展上下文窗口长度：更长的上下文意味着模型能“看到”更丰富的语境，从而做出更准确的预测。但这一切的前提是：必须使用高效注意力机制来让 $O(n^2)$ 的计算变得可承受。

\subsection{自注意力作为核回归：一种理论视角}

近年来，一些工作从核方法的角度对自注意力机制进行了解释，为理解其工作原理提供了优美的理论框架。Tsai 等人\cite{tsai2019transformerdissectionunifiedunderstanding} 指出，自注意力可以看作是对输入序列进行\textbf{非参数的核平滑}（non-parametric kernel smoothing）。具体而言，如果将 $QK^\top$ 中的内积 $\langle q_i, k_j\rangle$ 视为某种核函数 $k(q_i, k_j)$ 的评估值，那么 softmax 本质上就是对该核函数的归一化：

\[
\text{Attention}(Q, K, V)_i = \sum_{j=1}^n \frac{\exp(\langle q_i, k_j\rangle / \sqrt{d_k})}{\sum_{\ell=1}^n \exp(\langle q_i, k_\ell\rangle / \sqrt{d_k})} \cdot v_j.
\]

这正是 Nadaraya-Watson 核回归的形式：用核函数度量样本间的相似性，然后对对应的值进行加权预测。在这一视角下，“训练自注意力”相当于“学习一个好的核函数”，而 $W_Q$ 和 $W_K$ 正是定义核函数的参数。Wright 和 Gonzalez\cite{wright2021transformersdeepinfinitedimensionalnonmercer} 进一步证明，当 $W_Q$ 和 $W_K$ 的初始化满足一定条件时，自注意力层等价于可以被表征为一种在巴拿赫空间上的核学习方法。

这一理论视角的价值不仅在于提供了优雅的数学解释，更在于它指明了改进方向：如果自注意力本质上是核回归，那么我们是否可以设计更高效的核函数来替代标准的内积-softmax组合？事实上，Choromanski 等人\cite{choromanski2022rethinkingattentionperformers} 的 Performer 正是沿着这个思路，使用正随机正交特征（positive random orthogonal features）来近似 $\exp(q^\top k)$，从而避开了 $O(n^2)$ 的计算瓶颈。

\subsection{归纳偏置：为什么自注意力“学得快”？}

从统计学习理论的角度看，自注意力机制的成功也可以归因于其\textbf{归纳偏置}（inductive bias）与序列数据的天然结构高度匹配。Edelman 等人\cite{edelman2022inductivebiasesvariablecreation} 从理论上证明，有界范数的单头自注意力层可以表示输入序列的\textbf{稀疏函数}，即最终输出仅依赖于少数关键的输入位置。这对于语言等符号序列尤其重要：在一句话中，决定某个词语义的往往只是少数几个相关的词（例如“它的”所指代的名词），而非句子中的所有词。自注意力通过“查询-匹配-聚合”的机制，天然地倾向于学习这种稀疏的长距离依赖。

Yun 等人\cite{yun2020transformersuniversalapproximatorssequencetosequence} 则从\textbf{通用逼近定理}的角度出发，证明了 Transformer（自注意力加前馈网络）是序列到序列函数的通用逼近器。他们的关键结论是：只要有足够多的头和足够大的前馈网络，Transformer 可以在任意精度下逼近任何连续的序列到序列映射。这为自注意力机制的强大表达能力提供了理论基础。

这些理论分析共同指向一个核心事实：自注意力不仅仅是“能工作”，而是以\textbf{最少的先验结构假设}提供了\textbf{最大的表达灵活性}。正是这种“少假设、多学习”的设计哲学，使得 Transformer 架构能够在 NLP、计算机视觉、语音处理、蛋白质结构预测等截然不同的领域中都取得突破性进展\cite{bommasani2022opportunitiesrisksfoundationmodels}。

\subsection{为什么多头注意力能够捕捉复杂关系？}

上面的实验观察“多头能够自发分化出不同的功能偏好”并非偶然。近年来，大量研究从不同角度对这一问题进行了解释，虽然尚无统一的终极答案，但几条互补的理论线索已经浮现。

\paragraph{子空间多样性：每个头生活在不同的世界里}

最直接的解释来自多头注意力的数学结构本身。由于每个头拥有独立学习的一组投影矩阵 $W_Q^{(i)}, W_K^{(i)}, W_V^{(i)}$，每个头实际上在输入空间的不同子空间中计算注意力。即使所有头接收完全相同的输入 $X$，它们通过各自不同的线性投影将其映射到了不同的低维流形上。Voita 等人\cite{voita2019analyzingmultiheadselfattentionspecialized} 通过大规模探针实验对 Transformer 的注意力头进行了系统性分类，发现不同的头确实演化出了截然不同的功能角色：

\begin{itemize}
    \item \textbf{位置头}（positional heads）：注意力权重强烈集中于相邻位置（如前一个或后一个 token），近乎于学习了一个卷积核式的局部滑动窗口。这类头在底层尤为常见，负责提取局部短语结构。
    \item \textbf{语法头}（syntactic heads）：注意力集中于与当前 token 有特定句法依赖关系的 token（如动词的主语、名词的修饰形容词等），其行为与依存句法分析器高度一致。
    \item \textbf{语义头}（semantic heads）：注意力分布较为分散，倾向于关注与当前 token 语义相关但未必句法相邻的 token（如同义表达、共指关系等），在高层更为常见。
    \item \textbf{稀有词头}（rare-token heads）：对低频词或特殊符号（如分隔符、标点）有强烈偏好。
\end{itemize}

这种自发的功能分化在 Clark 等人\cite{clark2019doesbertlookat} 对 BERT 的注意力分析中也得到了印证。他们发现某些头在特定语言学现象（如介词宾语、限定词-名词关系）上的注意力模式与人类标注的句法结构高度吻合。重要的是，这些角色并非预设的，而是在端到端的训练中\textbf{自然涌现}的：模型发现将不同子空间分配给不同语言学功能，能使整体损失下降得更快。

\paragraph{集成学习的视角：多头即多专家}

另一种视角将多头注意力理解为一种\textbf{集成学习方法}。单头注意力在计算 $\text{softmax}(QK^\top/\sqrt{d_k})$ 时，通过 softmax 的指数性质，往往会将大部分概率质量集中在少数几个位置，从而形成一种相对“尖锐”的关注模式。当一句话中存在多种同时生效的关系时（例如一个词既是前面动词的宾语、又是后面关系从句的先行词），单头就难以兼顾。

Michel 等人\cite{michel2019sixteenheadsreallybetter} 对 Transformer 的注意力头进行了系统性的消融实验，他们的核心发现之一是：\textbf{在测试时，大量注意力头可以被移除而性能仅轻微下降}。在机器翻译任务上，甚至可以移除 48/96 个头而 BLEU 值仅下降不到 1 个点。但这并不意味着多头设计是多余的。恰恰相反，他们在训练动态中发现：不同的头在训练的不同阶段承担了不同的“探索”功能，使得优化过程能够从多个不同的初始化出发，避免陷入局部极小值。这类似于随机森林中的“特征子空间抽样”：每个头在不同的子空间中寻找有用的关联，最终通过 $W_O$ 矩阵将各头的信息进行加权融合。

这一点也得到了后续理论工作的支持。Hron 等人\cite{hron2020infiniteattentionnngpntk} 从高斯过程的角度证明，当头的数量趋于无穷且每个头的维度按比例缩小时，多头注意力收敛于一个与单头注意力具有不同核函数的极限过程。换句话说，多头结构\textbf{根本性地改变了模型的函数空间}，而不仅仅是增加了参数。

\section{NLP中的Transformer}
自然语言处理是 Transformer 架构的最初应用场景，也是直到现在 Transformer 架构最为成功的应用领域之一。其做法是：

\begin{enumerate}
    \item 对于输入的文本序列，经过tokenizer预处理，将其转换为一个\textbf{token序列}，每一个token是一个整数，表示词汇表中对应的词的索引。
    \item 对于每一个输入的token序列，首先将其转换为一个\textbf{词嵌入向量序列}，每一个词嵌入向量的维度为$d_{model}$。
    \item 将词嵌入向量序列输入到一个\textbf{编码器}中，编码器由若干个编码器层堆叠而成，每一层都包含一个多头自注意力层和一个FNN。
    \item 编码器的输出是一个上下文相关的词嵌入向量序列，每一个词嵌入向量的维度仍然为$d_{model}$。
    \item 将编码器的输出输入到一个\textbf{解码器}中，解码器由若干个解码器层堆叠而成，每一层都包含一个多头自注意力层、一个编码器-解码器注意力层和一个FNN。
    \item 解码器的输出是一个上下文相关的词嵌入向量序列，每一个词嵌入向量的维度仍然为$d_{model}$。
    \item 将解码器的输出输入到一个线性变换层和一个softmax层，得到一个词汇表大小的概率分布向量，每一个元素表示对应token的生成概率。
    \item 这个概率分布向量可以用于生成下一个token，或者用于计算损失函数进行模型训练。
\end{enumerate}

\subsection{Tokenizer}

Tokenizer是自然语言处理中的重要组件，作用就是将输入的文本序列转换为一个token序列。这个过程看似简单，实则蕴含着对“什么是语言的基本单位”这一根本问题的回答。tokenization 的质量直接影响模型对文本的建模能力：切分太粗，词汇表将极为庞大且充满稀疏项；切分太细，序列会变得过长，信息密度降低。现代 NLP 中，tokenizer 主要有三种类型：基于规则的、基于统计的、基于学习的。

\paragraph{基于规则的Tokenizer} 基于规则的Tokenizer是最早出现的类型，其主要思想是通过一些预定义的规则来将文本序列切分为token序列。例如，空格分词（按空白字符切割，英文中最朴素的做法）、标点符号分词（将标点独立成词）、正则表达式分词（用规则匹配数字、日期、URL等模式）等等。这一方法的优点在于简单、快速、透明，不需要任何训练数据或统计信息。然而，其缺点也同样明显：无法处理复杂的语言现象。对于形态丰富的语言（如芬兰语中一个动词可能有上百种变形），规则难以穷举；对于没有天然词间分隔符的语言（如中文、日文），空格规则直接失效；对于不断涌现的新词、拼写错误、社交媒体中的变体表达，固定规则更是无从下手。因此，纯规则分词在现代大模型时代已很少单独使用，更多是作为预处理中的一个辅助步骤。

\paragraph{基于统计的Tokenizer} 基于统计的Tokenizer是当前大语言模型中最主流的分词方案。其核心思想不再是“如何将已经定义好的词分开”，而是反过来：\textbf{从大量文本数据中自动发现哪些子字符串应该成为一个词}。其基本假设是：频繁共现且共现模式稳定的字符序列，很可能对应一个有意义的语言单位（词或词根）。目前被广泛使用的三种算法是 BPE、WordPiece 和 Unigram。

\textbf{字节对编码（Byte Pair Encoding, BPE）}最初是1994年提出的数据压缩算法\cite{gage1994new}，于2016年被 Sennrich 等人\cite{sennrich2016neuralmachinetranslationrare}引入到神经机器翻译中后大放异彩。其算法极其简洁：
\begin{enumerate}
    \item 将训练语料中的每个词拆成字符，并在词尾附加一个特殊符号（如“\verb|_|”或“\verb|</w>|”）以标记单词边界。
    \item 统计所有相邻符号对的出现频率。
    \item 将频率最高的符号对合并为一个新的符号，加入词汇表。
    \item 重复步骤2--3，直到词汇表达到预定大小，或没有可合并的符号对为止。
\end{enumerate}
BPE 的巧妙之处在于它自动发现了“次词”单元（subword units）：在英文中，常见的词根、词缀（如“ing”“tion”“un”）会被反复合并而进入词汇表，而罕见词会被切分成更小的子单元。这样一来，模型即使遇到训练时未见过的词（如一个新造的复合词），仍能通过其组成部分来推断意义，因为它至少见过“cat”和“like”，所以对“catlike”并不会完全陌生。

\textbf{WordPiece}\cite{6289079}是 Google 在 BERT 中采用的分词算法，其反向逻辑与 BPE 相似，但选择合并对的标准不是频率，而是最大似然：在所有候选符号对中，WordPiece 选择那个能使训练语料似然值（基于语言模型）增益最大的对进行合并。换言之，它倾向于合并那些“放在一起后使整个序列的概率增加最多”的字符对。这与 BPE 的纯频率驱动形成了差异：BPE 倾向于合并最常出现的对，而 WordPiece 倾向于合并最\textbf{不可分割}的对。例如，如果“s”和“ing”各自已经很好地建模了数据，即使“sing”频繁出现，WordPiece 也未必会将其合并。

\textbf{Unigram} 采用了一种不同的思路：它从一个初始大词汇表（例如包含所有可能的子串）开始，基于期望最大化算法反复移除那些使训练语料似然值下降最少的符号，最终收缩到目标词汇表大小。这与 BPE 从最小单元“向上合并”的过程正好相反，是一种“向下剪枝”的策略。SentencePiece\cite{kudo2018sentencepiecesimplelanguageindependent} 将 BPE 与 Unigram 统一在同一框架内，并特别地\textbf{将所有输入视为原始字节流，直接对字节进行分词}，从而彻底避免了因不同语言预分词规则不一致带来的工程麻烦。这是当前多语言大模型（如 LLaMA）的标配。

基于统计的tokenizer的核心优势在于：它通过大数据驱动的方式，自动发现了语言中稳定的“构词基元”，从而在\textbf{词汇表大小}和\textbf{序列长度}之间取得了极佳的折中。一个典型的大模型词汇表大小约为3万到10万（如 GPT-2 使用约5万的 BPE 词汇表，BERT 使用约3万的 WordPiece 词汇表），远小于自然语言的真实词汇量（通常以数百万计），又远大于字符集的规模（几十到几百）。这一折中极大地降低了 softmax 层的计算开销和嵌入矩阵的内存占用。

\paragraph{基于学习的Tokenizer} 基于学习的Tokenizer是近年来出现的一种新趋势，其核心思想是用一个可训练的神经网络来直接建模从字符串到token序列的映射。例如，可以用一个字符级别的 RNN 或 Transformer 来预测切分点，将 tokenization 本身作为下游任务的子模块进行端到端训练。这类方法的优点是能够内化特定任务对“好的分词”的定义。然而，其计算开销通常远高于统计方法，且在大规模预训练中尚未展现出对 BPE/Unigram 的显著优势，因此目前仍处于研究阶段，离工业界大规模采用还有距离。

\subsection{词嵌入向量}

有了 token 序列之后，计算机面临的第一个问题是：如何用数字来表示这些 token？计算机只能处理数值，而 token 是离散符号，“cat”和“dog”之间的抽象语义关系不会自动编码进整数索引中。我们需要一种将离散符号映射到连续向量空间的表示方法。

\paragraph{独热编码及其困境} 对于一个 token，最简单的表示方法是使用\textbf{独热编码}（One-hot Encoding）：将 token 在词汇表中的索引位置设为 1，其余位置设为 0。假设词汇表大小为 $V$，则每一个 token 的独热编码向量的维度为 $V$。在 $V = 50000$ 时，每个词被表示为一个长度为 50000 的向量，其中只有一个分量为 1，其余全为 0。这种表示有三个致命缺陷：

\begin{description}
    \item[维度灾难] 向量的维度等于词汇量，在典型的大模型场景中这意味着数万维的稀疏向量，存储和计算开销不可接受。
    \item[语义鸿沟] 任意两个不同 token 的独热向量之间的欧氏距离都是 $\sqrt{2}$，内积都是 0。“猫”和“狗”之间、“猫”和“电脑”之间在数学上毫无区别。模型无法从输入表示中感知到任何词与词之间的语义亲和性。
    \item[泛化缺失] 模型在训练中学习了“猫喜欢喝牛奶”，但对“狗喜欢喝牛奶”无从得知任何线索，因为“猫”和“狗”在输入空间中毫无关联。
\end{description}

\paragraph{分布式表示与词嵌入的思想} 解决上述困境的思路是\textbf{分布式表示}（distributed representation）：不应让每一个 token 占据一个正交的坐标轴，而应让它们“共享”一个低维连续空间的各个维度，每个 token 在这个空间中的位置编码了它的语义信息。在这一框架下，“猫”和“狗”的向量应当接近（因为它们同属宠物，语义相似），而与“电脑”的向量应当较远。更令人惊喜的是，这种向量还应当满足\textbf{语义加减}：例如，我们希望 $v_{\text{king}} - v_{\text{man}} + v_{\text{woman}} \approx v_{\text{queen}}$，即向量空间中能够编码“性别”这一语义维度，沿着这一方向移动可以改变词的性别属性而保持其“皇室”语义不变\cite{mikolov2013efficientestimationwordrepresentations,mikolov2013distributedrepresentationswordsphrases}。

这就是\textbf{词嵌入向量}（Word Embedding）：它将每个 token 映射到一个 $d_{\text{model}}$ 维的密集向量，这些向量在训练过程中会被学习到。在 Transformer 中，词嵌入是模型的第一层操作：设输入序列为 $(t_1, t_2, \dots, t_n)$，每个 $t_i \in \{1, \dots, V\}$，则嵌入层通过查表操作生成初始表示：
\[
\mathbf{x}_i = E_{t_i}, \quad E \in \mathbb{R}^{V \times d_{\text{model}}}.
\]
矩阵 $E$ 就构成了模型的词嵌入表，每一行是一个 token 的嵌入向量。在预训练过程中，$E$ 与模型的其它参数一起通过梯度下降进行优化。训练完成后，语义相似的 token 会发现它们的嵌入向量在高维空间中彼此靠近，而语义无关的 token 则彼此远离。

\paragraph{经典词嵌入方法及其与 Transformer 嵌入的关系} 在 Transformer 普及之前，词嵌入的研究已经经历了长足的发展，涌现出 Word2Vec\cite{mikolov2013efficientestimationwordrepresentations,mikolov2013distributedrepresentationswordsphrases}、GloVe\cite{2014Glove}、FastText\cite{bojanowski2017enrichingwordvectorssubword} 等方法。这些方法通过不同的训练目标来学习静态词嵌入：例如 Word2Vec 的 Skip-gram 模型通过预测上下文词来学习表示，GloVe 则将全局词共现矩阵分解为低维表示的结合。它们的共同特点是：为每一个 token 学习一个\textbf{固定的}嵌入向量，训练完毕后词的表示不再改变。

这与 Transformer 中的词嵌入有本质不同。在 Transformer 中，$E$ 矩阵在整个预训练过程中不断被更新，且它只是“初始表示”。经过自注意力和前馈网络的层层处理后，最终的\textbf{上下文表示}（contextualized representation）会根据周围词语动态调整。同一个词“bank”在“river bank”和“bank account”中，其初始嵌入是相同的，但在高层编码器输出中会产生截然不同的向量，因为自注意力机制将其上下文信息聚合进了表示中。这是静态词嵌入无法做到的事情，也是 Transformer 系列模型在 NLP 任务上全面超越传统方法的核心原因之一\cite{peters2018deepcontextualizedwordrepresentations,devlin2019bertpretrainingdeepbidirectional}。

\subsection{位置编码}

读者可能已经发现了一个致命漏洞：自注意力机制中的 $QK^\top$ 计算\textbf{不关心输入的顺序}。如果将输入序列的 token 全部打乱，注意力矩阵只会跟着重新排列，其数值不会改变，因为 $QK^\top$ 仅依赖于成对内积，而与 token 在序列中的绝对位置无关。这意味着，如果在输入中完全不加位置信息，那么句子“我打你”和“你打我”在模型看来将没有任何区别，这显然是错误的。

解决这一问题的方法是在信息流入 Transformer 层之前，向初始嵌入中注入\textbf{位置编码}（Positional Encoding），使得模型能够感知每个 token 的绝对或相对位置。

\paragraph{正余弦位置编码} 在原始 Transformer 论文\cite{vaswani2023attentionneed}中，Vaswani 等人提出了一种优雅的确定性位置编码——正余弦位置编码。对于位置 $pos$（从 $0$ 开始）和维度索引 $i$（从 $0$ 到 $d_{\text{model}}-1$），编码值为：
\begin{equation}
    PE_{(pos, 2i)} = \sin\left(\frac{pos}{10000^{2i / d_{\text{model}}}}\right), \quad
    PE_{(pos, 2i+1)} = \cos\left(\frac{pos}{10000^{2i / d_{\text{model}}}}\right).
\end{equation}
然后将位置编码直接加到词嵌入向量上作为模型的输入：
\[
\mathbf{x}_t^{\text{input}} = E_{t} + PE_t.
\]

这一设计的巧妙之处在于它同时满足了若干良好性质：
\begin{itemize}
    \item \textbf{值域有界}：正弦和余弦函数的值都在 $[-1, 1]$ 区间内，不会破坏输入的数值范围。
    \item \textbf{不同维度对应不同的频率}：$10000^{2i / d_{\text{model}}}$ 的指数项使得较低的维度对应较短的波长（高频，捕捉局部位置关系），较高的维度对应较长的波长（低频甚至远大于最大序列长度，捕捉全局位置关系）。这类似于二进制计数中不同位的进位频率不同。
    \item \textbf{相对位置的线性可表性}：对于固定的偏移 $k$，存在一个仅依赖于 $k$ 的线性变换（旋转矩阵），使得 $PE_{pos+k}$ 可以从 $PE_{pos}$ 通过线性变换得到。这意味着模型可以通过学习一个简单的线性映射来在不同位置之间迁移“相对偏移”的知识，而不需要为每一对相对距离独立学习参数。
\end{itemize}

从波的角度来看，正弦位置编码相当于为每个 token 分配了一个独特的“信号”，该信号在序列上的传播具有确定的相位关系。模型通过学习 $W_Q$ 和 $W_K$，可以将这些信号中的有用成分（例如“相邻位置有相似的编码”）提取出来，辅助注意力计算。

\paragraph{可学习位置编码} 另一种思路更加直接：为什么不让模型自己学会如何表示位置呢？可学习位置编码就是一个与嵌入矩阵同等地位的参数矩阵 $P \in \mathbb{R}^{\text{max\_len} \times d_{\text{model}}}$，其中第 $pos$ 行就是位置 $pos$ 的编码向量，与词嵌入相加后作为输入。BERT\cite{devlin2019bertpretrainingdeepbidirectional}、GPT\cite{2018Improving} 等早期预训练模型采用的就是这种方法。

可学习位置编码的优点是不需要手动设计任何数学函数，模型自动找到最优的位置表示方式。但其缺点也很明显：它天然地受限于训练时设定的最大序列长度。如果推理时遇到了比训练时更长的序列，对应的位置编码行就没有经过训练，模型行为不可预测（通常会直接失败或产生剧烈退化）。

\paragraph{相对位置编码} 正弦位置编码和可学习位置编码都属于\textbf{绝对位置编码}家族：每个位置拥有一个独立的编码，相加到输入表示中。但许多研究者认为，对于自注意力而言，更重要的是\textbf{相对位置}，即 token A 与 token B 之间的距离，而非它们在序列中的绝对坐标。这一思想催生了多种相对位置编码方案。

Shaw 等人\cite{shaw2018selfattentionrelativepositionrepresentations} 在 2018 年首次系统性地研究了相对位置编码在 Transformer 中的应用。他们修改了注意力的计算方式，在 $QK^\top$ 中直接加入一个表示相对距离的可学习偏置项：
\[
\text{Attention}(Q, K, V)_i = \sum_{j} \text{softmax}\left(\frac{q_i^\top k_j + a_{i-j}}{\sqrt{d_k}}\right) v_j,
\]
其中 $a_{i-j}$ 是一个基于相对距离的可学习参数。这意味着，模型不是通过“位置 i 的嵌入”来识别位置，而是通过“i 比 j 靠后几个位置”来调整注意力权重。这种方式更贴近语言的实际规律：一个词前的第2个词与它之间的关系，往往和这个词前的第50个词与它之间的关系有本质不同，而这种不同与绝对位置基本无关。

此后，Transformer-XL\cite{dai2019transformerxlattentivelanguagemodels} 进一步将相对位置编码与正弦编码相结合，实现了超长序列上的有效建模；T5\cite{raffel2023exploringlimitstransferlearning} 使用了简化的可学习相对位置偏置，将偏置值的数量封装在一个固定大小的“桶”中，减小了参数数量；DeBERTa\cite{he2021debertadecodingenhancedbertdisentangled} 解耦了内容注意力和位置注意力，分别计算 $QK^\top$ 的内容项和位置项，再用相对位置编码调制二者的交互。

相对位置编码已成为现代大语言模型的主流选择。例如，LLaMA\cite{touvron2023llamaopenefficientfoundation} 采用了源自 GPT-NeoX\cite{black2022gptneox20bopensourceautoregressivelanguage} 的旋转位置编码（RoPE）\cite{su2023roformerenhancedtransformerrotary}，而 PaLM\cite{chowdhery2022palmscalinglanguagemodeling} 则使用了 ALiBi\cite{press2022trainshorttestlong}。RoPE 通过将 Query 和 Key 在频域进行旋转，使得内积 $q_i^\top k_j$ 仅依赖于相对距离 $i-j$，同时保留了正弦位置编码的所有良好性质，在扩展序列长度时表现尤其优异。ALiBi 则更为激进地简化了位置编码：它完全取消了嵌入层的加法，改为在注意力分数中减去一个与相对距离成正比的偏置，让模型完全通过这一惩罚项来感知距离，在长距离外推上展现了令人意外的鲁棒性。

可以说，从绝对位置编码到相对位置编码，再到 RoPE 和 ALiBi 等精巧设计的出现，位置信息的注入方式已经成为一个丰富的研究子领域。这些设计的共同目标始终如一：\textbf{让模型以最高效、最泛化的方式理解“顺序”这一语言的根本结构}。没有位置编码，Transformer 不过是一个集合函数的逼近器；有了位置编码，它才能真正的处理人类语言。

\subsection{BERT与GPT：编码器与解码器的分化}

在前文所述的原始 Transformer 中，编码器和解码器是作为一个整体被设计出来处理序列到序列任务的：编码器负责“理解”源序列，将其转化为一组密集的上下文表示；解码器则以自回归的方式，依据这些表示和已经生成的部分目标序列，逐个词地产生输出。然而，当 Transformer 的强大能力被迁移到大规模预训练语料上时，研究者很快发现：\textbf{不一定需要完整的编码器-解码器结构，仅使用其中一部分就足以在大量任务上取得优异表现}。由此，形成了基于 Transformer 的两个最具影响力的预训练流派——以 BERT 为代表的\textbf{编码器模型}，和以 GPT 为代表的\textbf{解码器模型}。

\paragraph{BERT：双向编码器预训练}

BERT\cite{devlin2019bertpretrainingdeepbidirectional} 由 Google 在 2018 年提出，其核心思路极为明确：既然编码器中的自注意力机制天然是双向的（每个位置都能同时“看到”序列中所有其他位置），那么直接抽取一个多层 Transformer 编码器，在超大规模文本上进行无监督预训练，就能让模型学习到蕴含丰富上下文信息的文本表示。具体来说，BERT 就是一个由 $L$ 层 Transformer 编码器堆叠而成的网络，没有任何解码器组件。

与原始编码器-解码器架构的根本区别在于：BERT 不从输入本身自回归地生成输出，而是通过\textbf{破坏-重建}式的预训练任务来驱动学习。它使用了两种任务：
\begin{enumerate}
    \item \textbf{掩码语言模型（Masked Language Model, MLM）}：随机遮盖输入序列中 15\% 的 token，让模型预测这些被遮盖的 token。由于编码器在每个位置都可以看到除自己被遮盖部分以外的全部上下文（通过随机替换和保持原词的小技巧，迫使模型不能简单地“抄写”当前位置），模型必须学会利用双向信息进行推理。
    \item \textbf{下一句预测（Next Sentence Prediction, NSP）}：输入两个句子，判断它们在原文中是否相邻，促使模型学习句子间的关系。
\end{enumerate}
预训练完成后，BERT 就可以被用于各种自然语言理解下游任务：文本分类、句子对推理、命名实体识别、抽取式问答等。对于这些任务，只需在 BERT 的顶层输出上添加一个轻量的任务专用层（例如一个线性分类器），并用标注数据进行微调。由于 BERT 的输出是序列中每个 token 的上下文相关表示，它既可以直接用于 token 级别的标注任务，也可以提取首个 token 的特殊表示（[CLS]）用于句子级别的分类。

因此，BERT 与原始编码器-解码器架构的分野十分清晰：
\begin{itemize}
    \item \textbf{结构上}：BERT 只保留了 Transformer 的编码器，彻底舍弃了解码器和编码器-解码器注意力。这使得它不能进行序列生成，但结构更简单、更规整。
    \item \textbf{注意力模式上}：BERT 采用全双向自注意力，没有因果掩码，因此每个 token 都可以自由地聚合来自两边的信息，这对于理解任务至关重要。而原始编码器虽然也是双向的，但其解码器部分是单向（带掩码）的。
    \item \textbf{预训练目标上}：BERT 以“理解”为训练导向（填词、判断句子关系），而非“生成”。
\end{itemize}

\paragraph{GPT：因果解码器预训练与生成能力}

几乎在同一时期，OpenAI 推出了 GPT（Generative Pre-Training）\cite{2018Improving} 及其后续版本 GPT-2、GPT-3\cite{brown2020languagemodelsfewshotlearners}。GPT 走向了另一条路：它只使用 Transformer 的解码器，并赋予其大规模自回归语言建模的能力。

原始解码器在训练时就已经使用了因果掩码（causal mask）来确保对未来位置的不可见性，GPT 完全继承了这一点。一个由 $L$ 层 Transformer 解码器堆叠而成的网络，在预训练时只做一件事：给定前 $t-1$ 个 token，预测第 $t$ 个 token 的概率分布。这是一个典型的自回归语言模型目标：
\[
\mathcal{L} = \sum_{i} \log P(x_i \mid x_1, \dots, x_{i-1}; \Theta).
\]
与 BERT 的最大区别在于，GPT 没有使用任何“填空”或“破坏输入”的手段，而是直接从左到右地学习下一个词的分布。这使得它天然具备\textbf{文本生成}能力：预训练完成后，只需提供一个开头，模型就可以递归地采样下一个 token，不断续写出连贯的文本。

GPT 与原始编码器-解码器架构的不同体现在：
\begin{itemize}
    \item \textbf{结构上}：GPT 舍弃了编码器，仅保留解码器。值得注意的是，原始解码器内部包含自注意力层（带掩码）和编码器-解码器注意力层，而 GPT 的解码器中不存在后者，因为它没有一个单独的编码器来提供外部记忆。因此，GPT 的所有层都是仅依赖已生成文本的自注意力。
    \item \textbf{注意力模式上}：GPT 的注意力是严格单向的（自回归），不能像 BERT 那样同时看到两侧信息。这牺牲了一定的理解能力，但换取了任意长度文本的流畅生成能力。
    \item \textbf{预训练目标上}： “学会下一个词”。奇妙的是，正是这个看似简单的目标，在足够大的模型和数据的加持下，迫使模型内化了语法、事实知识、推理链条乃至一定程度的逻辑能力\cite{brown2020languagemodelsfewshotlearners}。
\end{itemize}

\paragraph{编码器-解码器预训练：保留序列到序列的灵活性}

尽管 BERT 和 GPT 分别在理解和生成上大放异彩，但有一类任务天然最适合完整的编码器-解码器结构：即那些输入和输出都是序列，且长度可能不同的任务，如机器翻译、文本摘要、对话生成等。对于这些任务，编码器需要先“消化”整个输入序列，生成一组上下文表示，然后解码器自回归地生成输出序列，并在每一步通过交叉注意力从中抽取相关信息。这就是原始 Transformer 的设计初衷。

因此，研究者也自然地探索了\textbf{在完整编码器-解码器框架上进行预训练}的路径。两个代表性工作是 T5（Text-to-Text Transfer Transformer）\cite{raffel2023exploringlimitstransferlearning} 和 BART（Bidirectional and Auto-Regressive Transformers）\cite{lewis2019bartdenoisingsequencetosequencepretraining}。它们的共同思想是：将几乎所有 NLP 任务统一表示为“文本到文本”的形式（输入一段文本，输出一段文本），然后使用编码器-解码器模型进行预训练。

T5 的编码器和解码器结构与原始 Transformer 类似，但预训练时采用了类似于 BERT 的 MLM 变体（遮盖连续片段，要求解码器生成被遮盖的内容），从而同时获得了双向编码器的理解能力和自回归解码器的生成能力。BART 则在编码器端引入双向注意力（类似 BERT），在解码器端保持自回归特性，训练目标是让模型学会重建被破坏的文本（例如打乱句子顺序、遮盖词汇等）。

与 BERT 和 GPT 相比，完整的编码器-解码器预训练模型具有以下特点：
\begin{itemize}
    \item \textbf{通用性最强}：既能做理解（编码器输入，在解码器上添加分类头），也能做生成（编码器输入，解码器自回归输出），并且天然适合序列到序列任务。
    \item \textbf{计算开销更大}：因为同时包含了编码器和解码器，参数量和训练时的计算量通常是同等层数的编码器模型或解码器模型的两倍左右。
    \item \textbf{灵活性}：编码器和解码器可以独立初始化、微调或与其他模块拼接，适用于复杂的多模态或多任务架构。
\end{itemize}

\paragraph{三种范式的对比总结}

\begin{table}
    \centering\small
    \begin{tabular}{c|p{4cm}p{4cm}p{4cm}}
        \toprule
        \textbf{特征} & \textbf{BERT（编码器）} & \textbf{GPT（解码器）} & \textbf{编码器-解码器} \\
        \midrule
        组成模块 & 仅编码器 & 仅解码器（无交叉注意力） & 编码器、解码器（含交叉注意力） \\
        注意力模式 & 全双向（无掩码）& 单向因果掩码 & 编码器双向，解码器因果掩码 + 交叉注意力 \\
        预训练目标 & 掩码语言模型 + 下一句预测 & 自回归语言模型 & 文本重建（如去噪、去遮盖） \\
        典型任务 & 文本分类、序列标注、自然语言推理 & 文本生成、对话、思维链推理 & 翻译、摘要、文本到文本统一任务 \\
        代表模型 & BERT, RoBERTa, DeBERTa & GPT, GPT-2, GPT-3, GPT-4 & T5, BART, mBART \\
        \bottomrule
    \end{tabular}
    \caption{三种基于 Transformer 的预训练架构对比}
    \label{tab:bert_gpt_comparison}
\end{table}

从表 \cref{tab:bert_gpt_comparison} 中可以看出，选择哪种架构本质上取决于\textbf{对“双向上下文”和“生成能力”之间权衡}的需求。BERT 因双向编码在理解任务上精度极高，但无法直接生成文本；GPT 凭借自回归生成能力可以续写任意长文，但每一时刻只能看见“过去”；编码器-解码器模型则兼具二者，但代价是更高的参数量和训练复杂度。在大模型时代，随着 GPT 系列通过思维链提示等方式在理解任务上逼近甚至超越了 BERT 系列，边界正在变得模糊，但这种结构性的差异依然是理解不同预训练范式作用的根本框架。

\paragraph{那为什么GPT“一统天下”了呢？} 读者可能会问：既然 BERT 在理解任务上表现优异，T5/BART 在序列到序列任务上也有优势，为什么在大模型时代，GPT 系列几乎成为了事实上的标准？这背后有几个原因。

我们知道，GPT的目标是“学会下一个词”，所以聊天机器人中只需要按以下方式使用：
\begin{enumerate}
    \item 用户输入一些内容。
    \item 通过其他脚本把这些内容转成 token 序列，然后过一遍词嵌入得到输入的表示矩阵 $X$ 。
    \item 由一个系统进行拼接：“\verb|[BOS]balabala<stop> [agent]|”（这只是一个示例）。
    \item GPT 把这个当成自己说过的话，然后以上述内容为输入，预测下一个 token 是什么，如此反复进行自回归生成，直到生成了一个\verb|<stop>| token 为止（也可能是其他的终止条件）。
    \item 重复执行以上过程，用户每次输入新消息，就把新的消息拼接在之前的对话历史后面，继续自回归生成。
\end{enumerate}

所以对于T5/BART这类编码器-解码器架构而言，它们的架构在对话场景下有几个致命缺陷：
\begin{itemize}
    \item 双向架构又贵又慢：用户消息先跑一遍完整编码器（昂贵），得到深层表示；然后解码器再自回归生成。每一步生成都需要编码器的输出固定在那里，但编码器的计算量已经是 GPT 的 1.5-2 倍。然后用户发送下一条消息，我们需要把整个对话历史和这个新消息再过一遍编码器……那这个效率就太低了。而 GPT 只需要把新消息拼在后面就行，KV 缓存是可以复用的。
    \item 时序性差：传统双向架构的输入是一个完整的文本序列，模型在训练时就习惯了这种完整的输入，所以倾向于把输入当成一个整体，输出也当成一个整体，于是在输出中容易出现与用户消息强相关但空泛的内容，以及“知道自己未来要写什么”导致的忽略时序问题（其实这个解释很感性，所以看来很牵强）；
    \item 训练目标和推理目标不一致：传统双向架构的目标是学会填空和判断两句话是否连续，这与我们的实际使用场景不一致。所以为了用这东西做对话，要么去掉双向的注意力（那就变成 GPT 了），要么保留双向但改用更复杂的训练目标（如 XLNet 的排列语言建模），但这东西会大幅增加其复杂度。
\end{itemize}

\section{Transformer 架构的跨领域影响}

Transformer 最初为机器翻译这一典型的序列到序列任务而设计，但其核心组件——自注意力机制——的通用性远超狭隘的语言序列。将自注意力理解为“从集合中动态选取相关信息”的通用算子，立刻打开了将其迁移到其他数据模态的大门。近年来，Transformer 已经在计算机视觉、语音处理、多模态学习乃至科学发现等领域取得了革命性进展，其影响之深远，使得“万物皆可 Transformer”从一句调侃逐渐变为现实。本节将简要回顾 Transformer 在几个典型领域中的核心思路与代表性工作。

\subsection{计算机视觉：从图像到“图像块序列”}

卷积神经网络（CNN）长期在计算机视觉领域占据统治地位，其平移等变性、局部感受野和多尺度层次结构天然匹配图像的二维拓扑结构。将 Transformer 应用于视觉的核心问题在于：\textbf{如何将连续的二维图像数据转化为自注意力能够处理的离散序列？}

Vision Transformer (ViT)\cite{dosovitskiy2021image} 提出了一个极为简洁而直接的方案：将图像分割成固定大小的块（patches，例如 $16\times16$），将每个块线性投影为一个向量，然后视为一个 token，连同位置编码一并送入标准的 Transformer 编码器。在足够规模的预训练数据（如 JFT-300M）上，ViT 无需任何卷积归纳偏置即可达到甚至超越当时最优 CNN 模型的分类精度。这一结果正式宣告：\textbf{自注意力可以天然地捕捉图像中的长程空间依赖}，打破了卷积核局限于邻域的束缚。

然而，ViT 直接处理图像块序列的 $O(n^2)$ 复杂度在输入分辨率稍高时便难以承受，且缺乏 CNN 天然具备的多尺度层次结构。Swin Transformer\cite{liu2021swin} 针对性地设计了移位窗口注意力（shifted window attention），将自注意力限制在不重叠的局部窗口内，并通过窗口的跨层移位实现窗口间信息交互，同时引入层次化下采样来构建类似 CNN 的特征金字塔。这种设计使 Swin Transformer 在目标检测、语义分割等密集预测任务上全面超越当时最优的 CNN 架构，并被广泛应用于计算机视觉的各个子领域。

与此同时，目标检测领域也出现了以 Transformer 为骨干的全新范式。DETR (Detection Transformer)\cite{carion2020endtoendobjectdetectiontransformers} 完全抛弃了经典的锚框、区域提议网络和非极大抑制等手工设计组件，将目标检测重塑为一个集合预测问题：用一个 CNN 骨干提取特征后，由 Transformer 编码器-解码器直接并行预测一组无序的边界框与类别标签，通过二分图匹配进行端到端训练。这种“大道至简”的思路开启了检测架构的新纪元，催生出后续 Deformable DETR\cite{zhu2021deformable} 等一系列高效变体。

\subsection{语音与音频处理：频谱图作为序列}

语音信号天然是时序序列，因此 RNN 曾是主流选择之一。然而，Transformer 的 $O(n^2)$ 自注意力对于长达数万帧的原始音频波形仍显昂贵。一个自然的折中是将音频信号转化为\textbf{梅尔频谱图}（mel spectrogram），频谱图在时间轴和频率轴上形成一个二维表示，既可以视作图像通过 ViT 方式处理（将频谱图切块），也可以沿时间轴视为一个长序列。

Audio Spectrogram Transformer (AST)\cite{gong2021ast} 正是借鉴 ViT 的思路，将音频梅尔频谱图划分为重叠的 $16\times16$ 时间-频率块，输入到标准的 Transformer 编码器中，在多个音频分类基准上取得了当时最优结果，证明了纯注意力机制在声学事件检测中的有效性。

在语音识别和翻译领域，OpenAI 的 Whisper\cite{radford2023robust} 是一个里程碑式的工作。Whisper 采用编码器-解码器 Transformer 架构，将输入的音频梅尔频谱图（编码器）直接映射为目标语言的文本 token（解码器），通过大规模弱监督训练实现了多语种识别、语音翻译和语言识别的统一。其编码器将音频谱图沿时间轴压缩后送入 Transformer，解码器自回归地生成文本，整个流程与文本到文本的 Transformer 几乎一致；唯一的变化是输入端模态从离散的 token 序列变成了连续的频谱图帧。这种架构上的高度统一正是 Transformer 泛化能力的生动体现。

\subsection{多模态学习：文本与图像的桥梁}

现实世界中的信息很少是单一模态的。一张图片往往伴随着标题，一段视频包含画面和声音。Transformer 天然的序列化抽象使得不同模态的数据可以“拼”成一个统一的序列，或者通过共享的注意力机制进行交互，从而催生了多模态大模型。

CLIP (Contrastive Language-Image Pre-training)\cite{radford2021learning} 是一个优雅而影响深远的架构：它包含一个图像编码器（ViT 或 CNN）和一个文本编码器（Transformer），通过对比学习迫使匹配的图文对在共享的嵌入空间中靠近，不匹配的远离。这一简单目标使得 CLIP 能够学习到强健的视觉-语言对齐，支持零样本图像分类（输入类别名称文本，直接与图像特征比对）和多种下游任务的迁移。此后，视觉问答、图像描述等任务纷纷借鉴这种双塔或融合编码器结构，将文本 token 与图像块嵌入共同输入 Transformer，由自注意力层跨模态地进行信息融合。

在图像生成领域，DALL·E\cite{ramesh2021zero} 将文本和图像 token 拼接为一个序列，使用自回归 Transformer 直接从文本和图像 token 的联合分布中生成图像。Stable Diffusion 等扩散模型中的去噪网络也基本采用 U-Net 结合 Transformer 注意力块的结构，利用交叉注意力将文本条件注入图像去噪过程。可以说，从理解到生成，Transformer 已成为多模态智能事实上的通用计算单元。

\subsection{科学发现与结构化数据}

Transformer 的影响力甚至超越了传统的感知领域，渗透到蛋白质结构预测、数学推导和药物发现等科学计算领域。

DeepMind 的 AlphaFold2\cite{jumper2021highly} 在蛋白质结构预测这一“50年难题”上取得了突破性成就，其核心架构 Evoformer 就是一个精巧设计的 Transformer 变体，能够在多序列比对（MSA）数据和残基对之间进行交叉注意力交互，迭代精修蛋白质的三维坐标。这里的“序列”不再是文字，而是氨基酸残基序列和进化信息。Transformer 的序列建模能力被巧妙地用于捕捉氨基酸之间的成对几何约束。

在数学和推理领域，最近的大语言模型（同样基于 Transformer 解码器架构）开始展现出初步的数学推导和定理证明能力；在表格数据、时序预测和时间序列分析中，Transformer 也逐渐取代传统的 RNN 和统计方法，通过自注意力捕获长程时间依赖。

这些应用反复印证了一个核心事实：Transformer 的成功并不在于它“理解”了语言，而在于它提供了一种\textbf{对任意标记集合间成对交互进行建模的通用框架}。只要数据能够被合理地序列化为 token，且 token 之间存在某种或局部或全局的相互依赖关系，Transformer 的自注意力机制就有可能从中抽取出有用的模式。这正是它能够在短短数年间从机器翻译走向千行百业、从文本走进多模态世界的根本原因。

